\begin{remarks}[7.9.10]
\label{IV.2.s7.9.10-fr}
(i) Mostraremos más adelante (18.8.11) que, para que un anillo local $B$ reducido
sea geométricamente unirrama, es necesario y suficiente que, para todo morfismo
étale $h:X'\to X=\operatorname{Spec}(B)$ y todo punto $x'\in X'$ sobre
el punto cerrado de $X$, $\mathcal O_{X',x'}$ sea íntegro. Se deduce de ello que, si
en (7.9.9) se supone no sólo que
$\operatorname{Spec}(B\otimes_Ak)$ sea geométricamente íntegro
en cada punto, sino que sea geométricamente unirrama, se puede concluir que $B$ es
\emph{geométricamente unirrama}.

\oldpage[IV-2]{223}

(ii) Sean $X$, $Y$ dos preesquemas excelentes y $f:X\to Y$ un morfismo
plano, y supongamos que, para todo morfismo finito $Y_1\to Y$, se pueda resolver
$(Y_1)_{\mathrm{red}}$. De (7.9.8), tomando como
$\boldsymbol R$ la propiedad de ser \emph{regular}, se deduce que el conjunto $U$ de los puntos
$x\in X$ tales que la fibra en el punto $f(x)$ del morfismo
\[
\operatorname{Spec}(\mathcal O_x\otimes_{\mathcal O_{f(x)}}\boldsymbol{k}(f(x)))
\longrightarrow\operatorname{Spec}(\boldsymbol{k}(f(x)))
\]
sea \emph{geométricamente regular}, es \emph{estable por generización}. No sabemos si
este conjunto es \emph{abierto} (o, lo que equivale a lo mismo (\textbf 0$_{\mathrm{III}}$,
9.2.5), si es \emph{construible}), incluso en el caso particular en que $Y$ sea el
espectro de $\mathbf Z$ o de un anillo de polinomios en una indeterminada $k[T]$
sobre un cuerpo $k$ y en que, por tanto, la condición de «posibilidad de resolución» impuesta a
$Y$ se satisfaga trivialmente.
\end{remarks}

\begin{center}
\emph{(Continuará.)}
\end{center}

\par\smallskip
{\small\noindent\textbf{Nota histórica de la edición.}
«No sabemos» reproduce la perspectiva de los autores en 1965; no afirma que el problema siga abierto actualmente.}\par
